\chapter{Quantum Algorithms for LP and SDP}
\label{ch:lec19}

\begin{objectives}
\begin{itemize}[leftmargin=1.4em,nosep]
  \item Recall the LP/SDP problems and their classical solvers.
  \item Explain the multiplicative-weights framework for SDP.
  \item Describe quantum SDP solvers based on Gibbs-state sampling.
  \item Outline quantum interior-point methods and their assumptions.
\end{itemize}
\end{objectives}

\todo{Write the introductory paragraph orienting the lecture.}

\section{Linear and semidefinite programming}
\label{sec:lec19-1}
\todo{Draft this section.}

\section{Multiplicative weights and the Arora--Kale framework}
\label{sec:lec19-2}
\todo{Draft this section.}

\section{Quantum SDP solvers via Gibbs sampling}
\label{sec:lec19-3}
\todo{Draft this section.}

\section{Quantum interior-point methods}
\label{sec:lec19-4}
\todo{Draft this section.}

\section{Speedups, caveats, and input/output costs}
\label{sec:lec19-5}
\todo{Draft this section.}

\exercisesection
\begin{exercise}
\todo{Add exercise.}
\end{exercise}
\begin{exercise}
\todo{Add exercise.}
\end{exercise}
\begin{exercise}
\todo{Add exercise.}
\end{exercise}

\readingsection
\begin{itemize}[leftmargin=1.4em]
  \item F.~Brand\\~ao, K.~Svore, \emph{Quantum Speed-ups for SDPs}, arXiv:1609.05537.
  \item J.~van Apeldoorn et al., \emph{Quantum SDP-Solvers}, arXiv:1705.01843.
  \item I.~Kerenidis, A.~Prakash, \emph{A Quantum Interior Point Method for LPs and SDPs}, arXiv:1808.09266.
\end{itemize}
